\chapter{Covering Spaces}
\label{ch:viii10-covering-spaces}

Covering spaces turn global topology into locally trivial sheets.  Around every
point of the base, a covering map looks like a disjoint union of copies of one
open neighborhood.  This simple local model gives strong global information:
fibers are discrete, connected covers have a constant number of sheets, the
circle admits the exponential cover by the real line, and coverings interact
directly with fundamental groups.

\section{Covering maps}

\begin{definition}[Covering map]
\label{def:viii10-covering-map}
A continuous surjection
\[
p:\widetilde X\to X
\]
is a covering map if every \(x\in X\) has an open neighborhood \(U\) such that
\[
p^{-1}(U)=\coprod_{\alpha\in A}V_\alpha
\]
with each \(V_\alpha\) open in \(\widetilde X\) and
\[
p|_{V_\alpha}:V_\alpha\to U
\]
a homeomorphism.
\end{definition}

Such a neighborhood \(U\) is called evenly covered.

\section{Sheets and fibers}

For \(x\in X\), the fiber is
\[
p^{-1}(x).
\]

\begin{proposition}[Fibers are discrete]
\label{prop:viii10-discrete-fibers}
Every fiber of a covering map is a discrete subspace of \(\widetilde X\).
\end{proposition}

If \(U\) is evenly covered, each sheet above \(U\) contains exactly one point
of the fiber over each \(x\in U\).

\section{Number of sheets}

\begin{proposition}[Constant sheet number on connected bases]
\label{prop:viii10-sheet-number}
If \(X\) is connected, then the cardinality
\[
|p^{-1}(x)|
\]
is independent of \(x\in X\).
\end{proposition}

A cover with \(n\) points in each fiber is called an \(n\)-sheeted cover.

\section{The universal example}

\begin{example}[The exponential cover]
\label{ex:viii10-circle-cover}
The map
\[
p:\mathbb R\to S^1,
\qquad
p(t)=e^{2\pi it},
\]
is a covering map.  Every sufficiently short open arc in \(S^1\) is evenly
covered by countably many translated intervals.
\end{example}

The fiber over \(1\in S^1\) is
\[
p^{-1}(1)=\mathbb Z.
\]

\section{Finite circle covers}

For \(n\ge1\),
\[
p_n:S^1\to S^1,\qquad p_n(z)=z^n
\]
is an \(n\)-sheeted covering map.

\begin{example}
The map \(z\mapsto z^2\) is a two-sheeted covering of the circle by itself.
\end{example}

\section{Trivial covers}

\begin{definition}[Trivial covering]
\label{def:viii10-trivial-cover}
For a discrete set \(F\), the projection
\[
X\times F\to X
\]
is a trivial covering.
\end{definition}

Every covering is locally of this form over an evenly covered neighborhood,
but need not be globally a product.

\section{Local homeomorphisms}

\begin{proposition}[Coverings are local homeomorphisms]
\label{prop:viii10-local-homeomorphism}
Every covering map is a local homeomorphism.
\end{proposition}

\begin{warning}[The converse fails]
\label{warn:viii10-local-converse}
A surjective local homeomorphism need not be a covering map without additional
global control on the sheets.
\end{warning}

\section{Restriction of a covering}

\begin{proposition}[Restriction]
\label{prop:viii10-restriction}
If \(p:\widetilde X\to X\) is a covering and \(A\subset X\), then the restricted
map
\[
p^{-1}(A)\to A
\]
is a covering with the subspace topologies.
\end{proposition}

\section{Pullback coverings}

Given \(p:\widetilde X\to X\) and a map \(f:Y\to X\), define
\[
f^*\widetilde X
=
\{(y,\tilde x)\in Y\times\widetilde X:
f(y)=p(\tilde x)\}.
\]

\begin{theorem}[Pullback covering]
\label{thm:viii10-pullback}
The projection
\[
f^*\widetilde X\to Y
\]
is a covering map.
\end{theorem}

This gives a systematic way to transport coverings along maps.

\section{Connected coverings}

A covering space \(\widetilde X\) may be disconnected even when \(X\) is
connected.  Its connected components themselves cover the base when the base
is connected and locally path connected.

\begin{proposition}[Components of a covering]
\label{prop:viii10-components}
For a connected locally path-connected base, every connected component of a
covering space maps onto the base and is itself a covering space.
\end{proposition}

\section{Universal covers}

\begin{definition}[Universal covering space]
\label{def:viii10-universal-cover}
A covering
\[
p:\widetilde X\to X
\]
is universal if \(\widetilde X\) is simply connected.
\end{definition}

The real line is the universal cover of \(S^1\).

\begin{example}
For \(n\ge2\),
\[
S^n\to\mathbb RP^n
\]
is a two-sheeted cover, and \(S^n\) is simply connected.  Thus it is the
universal cover of \(\mathbb RP^n\).
\end{example}

\section{Evenly covered neighborhoods}

If \(U\subset X\) is evenly covered, choosing a point
\[
\tilde x\in p^{-1}(x),\qquad x\in U,
\]
selects a unique sheet \(V\) above \(U\).  The inverse
\[
(p|_V)^{-1}:U\to V
\]
is a local section.

\section{Sections}

\begin{definition}[Section]
\label{def:viii10-section}
A section of \(p:\widetilde X\to X\) is a map
\[
s:X\to\widetilde X
\]
such that
\[
p\circ s=\operatorname{id}_X.
\]
\end{definition}

Every covering has local sections over evenly covered neighborhoods.  A global
section is much stronger and can force a connected cover to be trivial in many
situations.

\section{Covering isomorphisms}

\begin{definition}[Isomorphism of coverings]
\label{def:viii10-cover-isomorphism}
Two coverings \(p_i:\widetilde X_i\to X\) are isomorphic if there is a
homeomorphism
\[
h:\widetilde X_1\to\widetilde X_2
\]
such that
\[
p_2\circ h=p_1.
\]
\end{definition}

The base is held fixed.

\section{Coverings and fundamental groups}

A connected covering determines a subgroup of the base fundamental group.
If
\[
p:(\widetilde X,\tilde x_0)\to(X,x_0),
\]
then
\[
p_*:\pi_1(\widetilde X,\tilde x_0)\to\pi_1(X,x_0)
\]
is injective.

\begin{theorem}[Fundamental-group injection]
\label{thm:viii10-pi1-injection}
For a covering map, the induced homomorphism on fundamental groups is injective.
\end{theorem}

The proof uses lifting uniqueness, developed systematically in VIII/11.

\section{Index and number of sheets}

Under the standard hypotheses used in covering-space classification, a
connected covering corresponds to a subgroup
\[
H\le\pi_1(X,x_0),
\]
and the number of sheets is the index
\[
[\pi_1(X,x_0):H].
\]

\section{Universal cover and the trivial subgroup}

For a simply connected universal cover,
\[
\pi_1(\widetilde X)=0,
\]
so the corresponding subgroup of \(\pi_1(X)\) is trivial.

\section{Classification preview}

For path-connected, locally path-connected, semilocally simply connected
spaces, connected coverings are classified by conjugacy classes of subgroups of
\(\pi_1(X)\).  Pointed connected coverings correspond to actual subgroups.

This theorem relies on the lifting machinery of VIII/11 and the deck-group
structure of VIII/12.

\section{Why coverings matter}

Covering spaces convert loops into sheet changes.  A loop in the base lifts
after a starting point is chosen; whether its lift closes or ends at another
point of the fiber detects its fundamental-group behavior.  The next chapter
develops this lifting principle precisely.


\section{Exercises}

\begin{exercise}\label{exr:viii10-01}
Define an evenly covered neighborhood.
\end{exercise}
\begin{hint}
Write the preimage as disjoint sheets.

% VIII-PEDAGOGY-HINT-v1.1 exr:viii10-01 append
Require the sheets to be open and pairwise disjoint, and check that each restriction maps onto the same entire neighborhood \(U\).
\end{hint}
\begin{solution}
An open \(U\subset X\) is evenly covered if \(p^{-1}(U)=\coprod V_\alpha\) and every \(p|_{V_\alpha}:V_\alpha\to U\) is a homeomorphism.
\end{solution}

\begin{exercise}\label{exr:viii10-02}
Define a covering map.
\end{exercise}
\begin{hint}
Every base point needs an evenly covered neighborhood.

% VIII-PEDAGOGY-HINT-v1.1 exr:viii10-02 append
A local inverse near one point upstairs is not enough: one neighborhood downstairs must simultaneously account for its whole preimage.
\end{hint}
\begin{solution}
A continuous surjection locally decomposing into disjoint homeomorphic sheets.
\end{solution}

\begin{exercise}\label{exr:viii10-03}
What is a fiber?
\end{exercise}
\begin{hint}
Take the inverse image of one point.

% VIII-PEDAGOGY-HINT-v1.1 exr:viii10-03 append
Distinguish the fiber \(p^{-1}(x)\) from the preimage of a neighborhood, which is a union of open sheets.
\end{hint}
\begin{solution}
\(p^{-1}(x)\).
\end{solution}

\begin{exercise}\label{exr:viii10-04}
Why are covering fibers discrete?
\end{exercise}
\begin{hint}
Use one evenly covered neighborhood.

% VIII-PEDAGOGY-HINT-v1.1 exr:viii10-04 append
Intersect one open sheet with \(p^{-1}(x)\); the homeomorphism onto an evenly covered neighborhood makes this intersection a singleton.
\end{hint}
\begin{solution}
Each sheet contains one fiber point, isolated inside that sheet.
\end{solution}

\begin{exercise}\label{exr:viii10-05}
What is an \(n\)-sheeted cover?
\end{exercise}
\begin{hint}
Count fiber points.

% VIII-PEDAGOGY-HINT-v1.1 exr:viii10-05 append
Count points over a fixed basepoint, not connected components of the total space; a connected cover can have several sheets.
\end{hint}
\begin{solution}
A covering whose fibers have exactly \(n\) points.
\end{solution}

\begin{exercise}\label{exr:viii10-06}
When is sheet number constant?
\end{exercise}
\begin{hint}
Use connectedness of the base.

% VIII-PEDAGOGY-HINT-v1.1 exr:viii10-06 append
An evenly covered neighborhood identifies its fibers with the same sheet-index set. Thus fiber cardinality is locally constant before connectedness is used.
\end{hint}
\begin{solution}
For a connected base, fiber cardinality is constant.
\end{solution}

\begin{exercise}\label{exr:viii10-07}
Give the universal cover of \(S^1\).
\end{exercise}
\begin{hint}
Use the exponential map.

% VIII-PEDAGOGY-HINT-v1.1 exr:viii10-07 append
Choose an arc shorter than a full turn and lift its angle to disjoint intervals differing by integers; also check that the total space is simply connected.
\end{hint}
\begin{solution}
\(p:\mathbb R\to S^1,\ p(t)=e^{2\pi it}\).
\end{solution}

\begin{exercise}\label{exr:viii10-08}
What is the fiber over \(1\) in the exponential cover?
\end{exercise}
\begin{hint}
Solve \(e^{2\pi it}=1\).

% VIII-PEDAGOGY-HINT-v1.1 exr:viii10-08 append
Divide the angle by \(2\pi\) and use the periodicity of the complex exponential to identify every solution, including negative ones.
\end{hint}
\begin{solution}
\(\mathbb Z\).
\end{solution}

\begin{exercise}\label{exr:viii10-09}
Give an \(n\)-sheeted self-cover of \(S^1\).
\end{exercise}
\begin{hint}
Use the power map.

% VIII-PEDAGOGY-HINT-v1.1 exr:viii10-09 append
For \(n\ge1\), solve \(z^n=w\) on a small arc and separate its \(n\) continuous root branches.
\end{hint}
\begin{solution}
\(z\mapsto z^n\).
\end{solution}

\begin{exercise}\label{exr:viii10-10}
Define a trivial covering.
\end{exercise}
\begin{hint}
Use a discrete fiber.

% VIII-PEDAGOGY-HINT-v1.1 exr:viii10-10 append
The discreteness of \(F\) makes every \(X\times\{a\}\) an open sheet mapping homeomorphically onto \(X\).
\end{hint}
\begin{solution}
The projection \(X\times F\to X\) for discrete \(F\).
\end{solution}

\begin{exercise}\label{exr:viii10-11}
What local property does every covering have?
\end{exercise}
\begin{hint}
Restrict to one sheet.

% VIII-PEDAGOGY-HINT-v1.1 exr:viii10-11 append
Restrict \(p\) to the unique sheet containing a chosen point upstairs; that restriction supplies the required local homeomorphism.
\end{hint}
\begin{solution}
It is a local homeomorphism.
\end{solution}

\begin{exercise}\label{exr:viii10-12}
Does every surjective local homeomorphism automatically give a covering?
\end{exercise}
\begin{hint}
Recall the warning.

% VIII-PEDAGOGY-HINT-v1.1 exr:viii10-12 append
Test the projection \(\mathbb R\sqcup(0,\infty)\to\mathbb R\) given by inclusion on each summand: compare fiber cardinalities near zero.
\end{hint}
\begin{solution}
Not in complete generality; global sheet control can fail.
\end{solution}

\begin{exercise}\label{exr:viii10-13}
What happens when a covering is restricted to a subspace \(A\subset X\)?
\end{exercise}
\begin{hint}
Restrict domain to the inverse image.

% VIII-PEDAGOGY-HINT-v1.1 exr:viii10-13 append
Intersect an evenly covered neighborhood with \(A\) and its sheets with \(p^{-1}(A)\); use the subspace topologies on both sides.
\end{hint}
\begin{solution}
\(p^{-1}(A)\to A\) is again a covering.
\end{solution}

\begin{exercise}\label{exr:viii10-14}
Define the pullback covering set.
\end{exercise}
\begin{hint}
Impose \(f(y)=p(\tilde x)\).

% VIII-PEDAGOGY-HINT-v1.1 exr:viii10-14 append
Form a subspace of \(Y\times\widetilde X\), not an unrestricted product: both coordinates must represent the same point of \(X\).
\end{hint}
\begin{solution}
\(f^*\widetilde X=\{(y,\tilde x):f(y)=p(\tilde x)\}\).
\end{solution}

\begin{exercise}\label{exr:viii10-15}
What map is the pullback covering projection?
\end{exercise}
\begin{hint}
Project to \(Y\).

% VIII-PEDAGOGY-HINT-v1.1 exr:viii10-15 append
Over \(f^{-1}(U)\), each sheet of \(p\) supplies a graph \(y\mapsto(y,s(f(y)))\); use these graphs to verify the local covering model.
\end{hint}
\begin{solution}
\((y,\tilde x)\mapsto y\).
\end{solution}

\begin{exercise}\label{exr:viii10-16}
Define a universal cover.
\end{exercise}
\begin{hint}
Require the total space to be simply connected.

% VIII-PEDAGOGY-HINT-v1.1 exr:viii10-16 append
Include path connectedness in simple connectivity and keep the covering condition separate from the condition on its total space.
\end{hint}
\begin{solution}
A covering \(p:\widetilde X\to X\) with simply connected \(\widetilde X\).
\end{solution}

\begin{exercise}\label{exr:viii10-17}
What universal cover does \(\mathbb RP^n\) have for \(n\ge2\)?
\end{exercise}
\begin{hint}
Use the antipodal quotient.

% VIII-PEDAGOGY-HINT-v1.1 exr:viii10-17 append
Check both ingredients: the antipodal quotient is a double cover, and \(S^n\) is simply connected in the stated range \(n\ge2\).
\end{hint}
\begin{solution}
\(S^n\to\mathbb RP^n\).
\end{solution}

\begin{exercise}\label{exr:viii10-18}
Define a section of a covering.
\end{exercise}
\begin{hint}
It is a right inverse.

% VIII-PEDAGOGY-HINT-v1.1 exr:viii10-18 append
Type the maps carefully: \(s:X\to\widetilde X\), so the right-inverse equation is an equality of maps from \(X\) to itself.
\end{hint}
\begin{solution}
A map \(s:X\to\widetilde X\) with \(p\circ s=\operatorname{id}_X\).
\end{solution}

\begin{exercise}\label{exr:viii10-19}
Do coverings always have local sections?
\end{exercise}
\begin{hint}
Use an evenly covered neighborhood.

% VIII-PEDAGOGY-HINT-v1.1 exr:viii10-19 append
Choose a sheet \(V\) above \(U\); the inverse of \(p|_V\) is continuous and its composite with \(p\) is \(\operatorname{id}_U\).
\end{hint}
\begin{solution}
Yes.
\end{solution}

\begin{exercise}\label{exr:viii10-20}
Define an isomorphism of coverings over \(X\).
\end{exercise}
\begin{hint}
Use a homeomorphism commuting with projections.

% VIII-PEDAGOGY-HINT-v1.1 exr:viii10-20 append
The commuting equation holds over the fixed base, and \(h\) must be a homeomorphism, not merely a bijection between fibers.
\end{hint}
\begin{solution}
A homeomorphism \(h:\widetilde X_1\to\widetilde X_2\) with \(p_2h=p_1\).
\end{solution}

\begin{exercise}\label{exr:viii10-21}
What does a covering induce on \(\pi_1\)?
\end{exercise}
\begin{hint}
Use VIII/11 lifting uniqueness.

% VIII-PEDAGOGY-HINT-v1.1 exr:viii10-21 append
Start with an upstairs loop whose projection is null-homotopic; lift the based null-homotopy to show the original upstairs class vanishes.
\end{hint}
\begin{solution}
An injective homomorphism \(p_*:\pi_1(\widetilde X)\hookrightarrow\pi_1(X)\).
\end{solution}

\begin{exercise}\label{exr:viii10-22}
What subgroup corresponds to a universal cover?
\end{exercise}
\begin{hint}
The total fundamental group is trivial.

% VIII-PEDAGOGY-HINT-v1.1 exr:viii10-22 append
Compute the image of the upstairs fundamental group under \(p_*\); simple connectivity makes the source group trivial.
\end{hint}
\begin{solution}
The trivial subgroup.
\end{solution}

\begin{exercise}\label{exr:viii10-23}
How is sheet number related to subgroup index in the standard classification setting?
\end{exercise}
\begin{hint}
Count cosets.

% VIII-PEDAGOGY-HINT-v1.1 exr:viii10-23 append
Identify the fiber with cosets of the covering subgroup after choosing a basepoint and a left/right action convention; count those cosets.
\end{hint}
\begin{solution}
It equals \([\pi_1(X):H]\).
\end{solution}

\begin{exercise}\label{exr:viii10-24}
What is the bridge to VIII/11?
\end{exercise}
\begin{hint}
Loops and paths must be lifted.

% VIII-PEDAGOGY-HINT-v1.1 exr:viii10-24 append
For a path downstairs, specify one starting point upstairs, then ask separately whether a lift exists, is unique, and returns to its starting fiber point.
\end{hint}
\begin{solution}
VIII/11 proves path lifting, homotopy lifting, uniqueness, and the general lifting criterion.
\end{solution}


\section{Solved problem dossiers}

\begin{problem}[Exponential evenly covered arc]\label{prob:viii10-01}
Show a sufficiently short open arc \(U\subset S^1\) is evenly covered by \(p(t)=e^{2\pi it}\).
\end{problem}
\begin{solution}
Choose \(U\) avoiding one point so that it admits a continuous argument interval \((a,b)\) of length less than \(1\). Then \(p^{-1}(U)=\coprod_{k\in\mathbb Z}(a+k,b+k)\), and \(p\) is a homeomorphism from each interval to \(U\).
\end{solution}

\begin{problem}[Two-sheeted circle cover]\label{prob:viii10-02}
Show \(p(z)=z^2\) is a two-sheeted cover of \(S^1\).
\end{problem}
\begin{solution}
Every point \(w\in S^1\) has two square roots.  On a sufficiently short arc choose a continuous square-root branch; the two branches give disjoint sheets each mapped homeomorphically onto the arc.
\end{solution}

\begin{problem}[Degree-\(n\) circle cover]\label{prob:viii10-03}
Show \(z\mapsto z^n\) is \(n\)-sheeted.
\end{problem}
\begin{solution}
Each \(w=e^{i\theta}\) has the \(n\) roots \(e^{i(\theta+2\pi k)/n}\), and short arcs admit \(n\) disjoint continuous root branches.
\end{solution}

\begin{problem}[Trivial finite cover]\label{prob:viii10-04}
Describe the \(n\)-sheeted trivial covering of \(X\).
\end{problem}
\begin{solution}
Take \(X\times\{1,\dots,n\}\to X\).  Its \(n\) components are sheets globally homeomorphic to \(X\).
\end{solution}

\begin{problem}[Discrete fiber proof]\label{prob:viii10-05}
Prove directly that \(p^{-1}(x)\) is discrete.
\end{problem}
\begin{solution}
Choose evenly covered \(U\ni x\).  Each fiber point lies in a different sheet \(V_\alpha\), and \(V_\alpha\cap p^{-1}(x)\) is that single point.
\end{solution}

\begin{problem}[Sheet-number constancy]\label{prob:viii10-06}
Why is fiber cardinality locally constant on the base?
\end{problem}
\begin{solution}
Over an evenly covered neighborhood, every fiber has one point in each sheet. Hence all fibers over that neighborhood have the same cardinality. On a connected base a locally constant cardinal-valued function is constant.
\end{solution}

\begin{problem}[Restriction cover]\label{prob:viii10-07}
Prove \(p^{-1}(A)\to A\) is a covering for arbitrary \(A\subset X\).
\end{problem}
\begin{solution}
If \(U\) evenly covers \(x\in A\), then \(U\cap A\) is open in \(A\), and the sheets \(V_\alpha\cap p^{-1}(A)\) map homeomorphically onto \(U\cap A\).
\end{solution}

\begin{problem}[Pullback cover]\label{prob:viii10-08}
Prove the pullback of a covering is a covering.
\end{problem}
\begin{solution}
Choose an evenly covered \(U\) around \(f(y)\). Over \(f^{-1}(U)\), the pullback splits into copies indexed by the sheets above \(U\), with each copy projected homeomorphically to \(f^{-1}(U)\).
\end{solution}

\begin{problem}[Projective cover]\label{prob:viii10-09}
Why is \(S^n\to\mathbb RP^n\) two-sheeted?
\end{problem}
\begin{solution}
Each projective point is a line with exactly two unit representatives \(x,-x\). Small neighborhoods of a line lift to disjoint neighborhoods around these two representatives.
\end{solution}

\begin{problem}[Universal projective cover]\label{prob:viii10-10}
For \(n\ge2\), why is \(S^n\to\mathbb RP^n\) universal?
\end{problem}
\begin{solution}
The antipodal quotient is a covering, and \(S^n\) is simply connected for \(n\ge2\).
\end{solution}

\begin{problem}[Local sections]\label{prob:viii10-11}
Construct a local section over an evenly covered neighborhood.
\end{problem}
\begin{solution}
Choose one sheet \(V\) over \(U\) and set \(s=(p|_V)^{-1}:U\to V\subset\widetilde X\).
\end{solution}

\begin{problem}[Global section warning]\label{prob:viii10-12}
Why can the exponential cover \(\mathbb R\to S^1\) have no global section?
\end{problem}
\begin{solution}
A global section would induce a homomorphism splitting \(p_*:0=\pi_1(\mathbb R)\to\pi_1(S^1)\cong\mathbb Z\), impossible because \(p\circ s=\operatorname{id}\) would imply \(p_*s_*=\operatorname{id}_{\mathbb Z}\).
\end{solution}

\begin{problem}[Covering isomorphism]\label{prob:viii10-13}
Show an isomorphism of coverings maps each fiber bijectively to the corresponding fiber.
\end{problem}
\begin{solution}
If \(p_2h=p_1\), then \(p_1(\tilde x)=x\) implies \(p_2(h(\tilde x))=x\).  Since \(h\) is a homeomorphism, the induced fiber map is bijective.
\end{solution}

\begin{problem}[Connected component cover]\label{prob:viii10-14}
Why does a component of a cover over a connected locally path-connected base map onto the base?
\end{problem}
\begin{solution}
The image of a component is open by local sheets and closed under local continuation; equivalently path lifting from one image point reaches points along base paths. Connectedness forces the image to be all of the base.
\end{solution}

\begin{problem}[Universal cover subgroup]\label{prob:viii10-15}
What subgroup of \(\pi_1(S^1)\) corresponds to \(\mathbb R\to S^1\)?
\end{problem}
\begin{solution}
The total space is simply connected, so the subgroup \(p_*\pi_1(\mathbb R)\) is \(0\le\mathbb Z\).
\end{solution}

\begin{problem}[Finite circle subgroup]\label{prob:viii10-16}
What subgroup of \(\pi_1(S^1)\cong\mathbb Z\) corresponds to \(z\mapsto z^n\)?
\end{problem}
\begin{solution}
The induced map on \(\pi_1\) is multiplication by \(n\), so the subgroup is \(n\mathbb Z\), of index \(n\).
\end{solution}

\begin{problem}[Product pullback]\label{prob:viii10-17}
Pull back the trivial cover \(X\times F\to X\) along \(f:Y\to X\). Identify the result.
\end{problem}
\begin{solution}
The pullback consists of triples effectively determined by \((y,a)\in Y\times F\), so it is naturally the trivial cover \(Y\times F\to Y\).
\end{solution}

\begin{problem}[Cover of a simply connected base]\label{prob:viii10-18}
What can be said about a connected covering of a simply connected, suitably nice base?
\end{problem}
\begin{solution}
The corresponding subgroup of the trivial fundamental group is trivial, so the connected covering is isomorphic to the identity cover under the standard classification hypotheses.
\end{solution}

\begin{problem}[Cover and local topology]\label{prob:viii10-19}
Why do base and total space have the same local topological type under a covering?
\end{problem}
\begin{solution}
Near every total point, choose the sheet over an evenly covered neighborhood; the covering restriction is a homeomorphism onto the base neighborhood.
\end{solution}

\begin{problem}[Fiber cardinality of disconnected cover]\label{prob:viii10-20}
Can a covering over a connected base be disconnected while having constant fiber size?
\end{problem}
\begin{solution}
Yes.  A trivial \(n\)-sheeted covering \(X\times\{1,\dots,n\}\to X\) has \(n\) disconnected components when \(X\) is connected.
\end{solution}

\begin{problem}[Identity cover]\label{prob:viii10-21}
Show \(\operatorname{id}_X:X\to X\) is a one-sheeted covering.
\end{problem}
\begin{solution}
Every open neighborhood is evenly covered by itself, with one sheet.
\end{solution}

\begin{problem}[Composition preview]\label{prob:viii10-22}
Under suitable hypotheses, why is a composition of covering maps again locally a covering?
\end{problem}
\begin{solution}
Choose evenly covered neighborhoods successively so the upper sheets map homeomorphically through the intermediate cover to the base neighborhood.  The local sheet decomposition composes.
\end{solution}

\begin{problem}[Classification preview]\label{prob:viii10-23}
State the subgroup object associated with a pointed connected cover.
\end{problem}
\begin{solution}
Choose \(\tilde x_0\) above \(x_0\); the associated subgroup is \(p_*\pi_1(\widetilde X,\tilde x_0)\le\pi_1(X,x_0)\).
\end{solution}

\begin{problem}[Bridge to lifting]\label{prob:viii10-24}
Why is lifting the mechanism behind covering classification?
\end{problem}
\begin{solution}
Loops lift uniquely from a chosen sheet.  Whether a lifted loop closes detects membership in the associated subgroup; maps lift precisely when their fundamental-group images land in that subgroup.
\end{solution}
